% Solutions of 22-greedy-proofs.
\section*{Chapter \ref{ch:b1:greedy-proofs} --- Greedy Algorithms and Their Proofs}

\begin{solution}{exo:b1:greedy-proofs:1}
By finishing time: $[1, 3)$ kept; $[2, 5)$ overlaps it; $[4, 7)$ kept; $[6, 8)$ and $[3, 9)$ overlap it; $[8, 10)$ kept;
$[7, 11)$ and $[9, 12)$ overlap. Three talks. No four fit, since the rule is optimal
(\cref{thm:b1:greedy-proofs:earliest-finish}).
\end{solution}

\begin{solution}{exo:b1:greedy-proofs:2}
Merges: $1 + 1 = 2$, $2 + 2 = 4$, $3 + 4 = 7$, $5 + 7 = 12$, $8 + 12 = 20$. Lengths 5, 5, 4, 3, 2, 1 for the frequencies
1, 1, 2, 3, 5, 8; \omterm{def:b1:greedy-proofs:prefix-code}{cost} $5 + 5 + 8 + 9 + 10 + 8 = 45$.
\end{solution}

\begin{solution}{exo:b1:greedy-proofs:3}
$1111 \mid 110 \mid 01 \mid 10 \mid 1110$: \code{notes}.
\end{solution}

\begin{solution}{exo:b1:greedy-proofs:4}
$289 = 200 + 50 + 20 + 10 + 5 + 2 + 2$: seven coins. For 10 with coins of 1, 3 and 4, the greedy gives $4 + 4 + 1 + 1$,
four coins, against $4 + 3 + 3$.
\end{solution}

\begin{solution}{exo:b1:greedy-proofs:5}
The long talk $[0, 100)$ starts first, is kept, and overlaps all the others: one talk instead of four. The exchange step
replaces the optimum's first talk by the greedy's; it needs the greedy's talk to finish no later, and a talk that starts
first may finish last, so the rest of the optimum no longer fits after it.
\end{solution}

\begin{solution}{exo:b1:greedy-proofs:6}
All weights are equal, and the two tie rules merge the symbols in opposite orders: the receiver's tree is the mirror of
the sender's. The bits $01\,00\,11$ of \code{bad} read \code{c}, \code{d}, \code{a} with the receiver's code. Safe ways:
fix the tie rule as part of the file format, so that both build the same tree; or send the code itself rather than the
frequencies, for instance the codeword lengths of a canonical code (\cref{exo:b1:greedy-proofs:8}).
\end{solution}

\begin{solution}{exo:b1:greedy-proofs:7}
\omcode{code/bachelor-1/22-greedy-proofs/exercises/platforms.py}{5}{15}{Platforms for a timetable}
This is the \omterm{def:b1:greedy-proofs:intervals}{maximum overlap} of closed intervals: walk the arrivals in order, first freeing the platforms of the trains
that left strictly before (a train leaving at the very minute another arrives still holds its platform). Sorting costs
$O(n \log n)$ and the walk $O(n)$. The test compares with the largest number of trains present at each minute.
\end{solution}

\begin{solution}{exo:b1:greedy-proofs:8}
\omcode{code/bachelor-1/22-greedy-proofs/exercises/canonical.py}{5}{13}{The canonical code}
Read as binary numbers, the codewords of one length are consecutive, and the first codeword of a longer length extends,
with zeros, the number just after the last shorter codeword: no shorter codeword can be its beginning. The lengths alone
determine the code, so the sender transmits $n$ small numbers, and both sides build the same codewords; the test checks
the lengths and the prefix property on codes built from Huffman lengths.
\end{solution}

\begin{solution}{exo:b1:greedy-proofs:9}
Let $L$ be the largest length. A codeword of length $\ell$ is the beginning of $2^{L - \ell}$ words of length $L$; for
two codewords these sets are disjoint, otherwise one codeword would begin the other. So $\sum_i 2^{L - \ell_i} \leq
2^L$. If every inner node has two children, replace two sibling leaves at depth $d$ by their parent: $2 \cdot 2^{-d} =
2^{-(d - 1)}$, the sum does not change; repeating down to the root alone gives 1.
\end{solution}

\begin{solution}{exo:b1:greedy-proofs:10}
In a way of paying with fewest coins, each coin $c^i$ with $i < k$ appears at most $c - 1$ times, since $c$ of them could
be replaced by one $c^{i+1}$. With $n_i < c$ for $i < k$, the sum $n_0 + n_1 c + \dots + n_{k-1} c^{k-1}$ is less than
$c^k$, so $n_k$ is the quotient of the amount by $c^k$, and the $n_i$ below are the digits in base $c$ of the remainder:
there is a single such way of paying. The greedy rule takes as many $c^k$ as fit, then as many $c^{k-1}$ as fit in the
rest, and so on, which computes exactly this quotient and these digits. The test checks $c = 2, 3, 5, 10$ up to 600.
\end{solution}

\begin{solution}{exo:b1:greedy-proofs:11}
If a task $i$ runs just before a task $j$ with $d_j < d_i$, exchange them. The tasks before and after are not moved;
$j$ finishes earlier; $i$ finishes when $j$ finished, at time $T$, and its lateness $T - d_i$ is less than $T - d_j$,
the lateness of $j$ before the exchange. So the largest lateness does not increase. Starting from an optimal order and
exchanging such neighbours until none is left reaches the order by deadline (equal deadlines can be put in any order,
which does not change the largest lateness) without increasing the largest lateness. \code{exercises/lateness.py} checks
it against all orders of up to six tasks.
\end{solution}

\begin{solution}{exo:b1:greedy-proofs:12}
For $i = 1$, $g_1$ finishes first of all talks. If $f(g_{i-1}) \leq f(o_{i-1})$, then $o_i$ starts at or after
$f(o_{i-1}) \geq f(g_{i-1})$, so it is compatible with $g_1, \dots, g_{i-1}$, which all finish by $f(g_{i-1})$; the
greedy keeps the first compatible talk in finishing order, so $f(g_i) \leq f(o_i)$. If the optimum had a talk
$o_{k+1}$ beyond the greedy's last $g_k$, it would start after $f(o_k) \geq f(g_k)$, be compatible with all the
greedy's talks, and the greedy would have kept a talk after $g_k$. So both have the same number of talks.
\end{solution}

\begin{solution}{pb:b1:greedy-proofs:1}
\textbf{1.} 1437 characters and 24 symbols. The ends of lines separate the titles; without coding them, the catalogue
could not be rebuilt.

\textbf{2.} 5 bits: $2^4 = 16 < 24 \leq 32 = 2^5$.

\textbf{3.} $5 \times 1437 = 7185$ bits.

\textbf{4.} \code{e} (195) and the space (182); \code{w} (11) and \code{f} (12).

\textbf{5.} Yes: two different words of the same length never begin one another.

\textbf{6.} With $\mathtt{e} = 0$, $\mathtt{t} = 01$, $\mathtt{h} = 1$, the bits $01$ are \code{t} or \code{eh}.

\textbf{7.} \code{w} (11) and \code{f} (12), into a tree of weight 23.

\textbf{8.} \code{g}: the leaves are numbered in the order of the symbols, and \code{g} comes before \code{u}; the queue
compares the numbers when the weights are equal. The second merge joins \code{v} (16) and \code{g} (17).

\textbf{9.} 23: each merge removes one tree, from 24 down to 1.

\textbf{10.} $24 - k$.

\textbf{11.} $n - 1$ merges of three queue operations, $O(\log n)$ each: $O(n \log n)$.

\textbf{12.} The proof of \cref{thm:b1:greedy-proofs:huffman} holds for any choice between equal weights: every such
choice gives an optimal tree, so all have the same \omterm{def:b1:greedy-proofs:prefix-code}{cost}, though the trees and their codewords may differ.

\textbf{13.} Length 3: \code{e}, space. 4: \code{t}, \code{h}, \code{o}, \code{r}, end of line, \code{a}, \code{n}. 5:
\code{d}, \code{l}, \code{s}, \code{i}, \code{c}, \code{y}. 6: \code{b}, \code{m}, \code{p}, \code{k}, \code{g},
\code{u}, \code{v}. 7: \code{f}, \code{w}.

\textbf{14.} $\frac{2}{8} + \frac{7}{16} + \frac{6}{32} + \frac{7}{64} + \frac{2}{128} = \frac{32 + 56 + 24 + 14 +
2}{128} = 1$.

\textbf{15.} $3 \cdot 377 + 4 \cdot 621 + 5 \cdot 279 + 6 \cdot 137 + 7 \cdot 23 = 1131 + 2484 + 1395 + 822 + 161 = 5993$,
as printed.

\textbf{16.} $7185 - 5993 = 1192$ bits, 16.6\,\%.

\textbf{17.} $1111 \mid 1101 \mid 101$: \code{t}, \code{h}, \code{e}.

\textbf{18.} When the bits read so far form a codeword, no other codeword begins with them: the symbol is decided
without looking at the next bit.

\textbf{19.} \code{f} and \code{w}, codewords $1100101$ and $1100100$, at depth 7, the largest: the two least frequent
symbols, as the lemma allows.

\textbf{20.} The only child could take the place of its parent, which shortens every codeword below by one bit and
lowers the \omterm{def:b1:greedy-proofs:prefix-code}{cost}.

\textbf{21.} $x$ and $y$ sit one level below where $z$ was, and $f_z = f_x + f_y$: $f_x(\ell_z + 1) + f_y(\ell_z + 1) =
f_z \ell_z + f_x + f_y$, the other leaves being unchanged.

\textbf{22.} For 1 symbol, the single \omterm{def:b1:induction-and-termination:rooted-tree}{leaf} is optimal. For $n + 1$ symbols, \omterm{met:b1:greedy-proofs:huffman}{Huffman's algorithm} merges the two lightest
$x$, $y$ and continues as \omterm{met:b1:greedy-proofs:huffman}{Huffman's algorithm} on $n$ symbols; by induction it builds an optimal tree $T'$ for them, and
\cref{lem:b1:greedy-proofs:merge} makes the final tree optimal.

\textbf{23.} The theorem compares codes that give a codeword to each character. A code for pairs of characters is a code
for another alphabet, of pairs, where the theorem applies again.

\textbf{24.} \emph{Huffman's theorem on the catalogue}: no \omterm{def:b1:greedy-proofs:prefix-code}{prefix code} for its 24 symbols encodes the catalogue in fewer
than 5993 bits, the \omterm{def:b1:greedy-proofs:prefix-code}{cost} of the program's \omterm{met:b1:greedy-proofs:huffman}{Huffman code}, against 7185 bits for the fixed-length code.
\end{solution}
